\documentclass[11pt,letterpaper]{article}
\PassOptionsToPackage{hyphens}{url}
\usepackage[margin=.83in,top=.8in,bottom=.8in,headheight=15pt]{geometry}
\usepackage[T1]{fontenc}
\usepackage{lmodern,amsmath,amssymb,mathtools,bm,booktabs,array,graphicx,tikz}
\usetikzlibrary{arrows.meta,positioning,calc,decorations.pathreplacing,patterns}
\definecolor{ink}{HTML}{17344D}\definecolor{accent}{HTML}{087E8B}\definecolor{pale}{HTML}{EDF5F5}\definecolor{gold}{HTML}{AD6B19}
\providecommand{\dd}{\,\mathrm{d}}\providecommand{\R}{\mathbb{R}}
\providecommand{\E}{\mathbb{E}}\providecommand{\Var}{\operatorname{Var}}
\providecommand{\Cov}{\operatorname{Cov}}\providecommand{\diag}{\operatorname{diag}}
\providecommand{\norm}[1]{\left\lVert#1\right\rVert}
\providecommand{\abs}[1]{\left|#1\right|}
\usepackage{microtype,hyperref,fancyhdr,enumitem,amsthm,sectsty,needspace}
\allsectionsfont{\raggedright}
\hypersetup{colorlinks=true,linkcolor=ink,urlcolor=accent,pdfauthor={Zan Ahmad},pdftitle={Membrane currents and Hodgkin--Huxley gate kinetics / Review}}
\setlength{\parindent}{0pt}\setlength{\parskip}{6pt}\linespread{1.07}
\setlist{itemsep=3pt,topsep=4pt,leftmargin=*}
\emergencystretch=2em\clubpenalty=10000\widowpenalty=10000
\pagestyle{fancy}\fancyhf{}\fancyhead[L]{\small Mathematical Modeling}\fancyhead[R]{\small Lecture 9}
\fancyfoot[L]{\footnotesize Zan Ahmad}\fancyfoot[R]{\footnotesize\thepage}\renewcommand{\headrulewidth}{.3pt}
\newtheorem{proposition}{Proposition}\newtheorem{theorem}{Theorem}\newtheorem{lemma}{Lemma}
\theoremstyle{definition}\newtheorem{definition}{Definition}\newtheorem{example}{Example}
\numberwithin{equation}{section}
\begin{document}
\section*{Membrane currents and Hodgkin--Huxley gate kinetics --- review sheet}
Use this after attempting the worked examples and practice. Explain each claim in words before calculating.

\section*{Key ideas}
\begin{itemize}
\item Current balance controls voltage rate under a declared sign convention.
\item An ion current needs both conductance and a driving difference.
\item m, h and n are fractions; their combinations determine conductances.
\item A voltage clamp fixes rates; a free spike couples gates and voltage.
\item Equal pulse timing and initial conditions support a controlled comparison.
\item Numerical accuracy and biological validity require different evidence.
\end{itemize}
\section*{Common errors and repairs}
\textbf{Treating h as sodium current.} h is a dimensionless availability fraction; current also needs m, a maximal conductance and V - E\_Na.

\textbf{Calling a clamped-voltage gate trace a predicted voltage spike.} The clamp supplies the voltage. A free-voltage experiment must solve current balance.

\textbf{Assuming bounded gates make a coarse trajectory accurate.} Bounds check admissibility. Compare smaller steps and independent references for accuracy.

\section*{Retrieve without your notes}
\begin{itemize}
\item What is current through a pathway at its reversal potential?
\item For alpha =\allowbreak{} beta =\allowbreak{} 1/ms, what is the equilibrium gate fraction?
\item Does doubling membrane area at fixed density double the voltage input per area?
\end{itemize}
\section*{Readiness check}
\begin{itemize}
\item Use a stated current sign convention to predict whether membrane voltage rises or falls.
\item Interpret m, h and n as gate fractions and connect them to conductance.
\item Distinguish a voltage clamp from a freely evolving voltage.
\item Compare two controlled stimuli and identify positive feedback and recovery.
\item Separate numerical checks from evidence about an actual neuron or heart cell.
\end{itemize}
\end{document}
